% ══════════════════════════════════════════════════════════════════════
% Chapter 3: Costly Depth and Ritual Sacrifice
% ══════════════════════════════════════════════════════════════════════

\chapter{Costly Depth and Ritual Sacrifice}
\label{ch:costly}

\epigraph{``A signal that is cheap to produce is a signal that is easy
to fake.''}{Amotz Zahavi, \textit{The Handicap Principle}}

% ── Introduction ──────────────────────────────────────────────────────

Why do organisms---and human societies in particular---invest enormous
resources in signals that appear, on the surface, to be wasteful?  The
peacock drags a metabolically expensive tail through the underbrush; the
stotting gazelle leaps conspicuously in the presence of a predator
rather than fleeing at full speed; the young Maasai warrior hunts a lion
with a spear when safer alternatives exist.  In each case, the very
costliness of the signal is what makes it \emph{honest}: a weak peacock
cannot afford the tail, a slow gazelle cannot afford the leap, and a
cowardly warrior cannot afford the hunt.\footnote{Zahavi's original
  formulation (1975) was widely resisted until Grafen (1990) provided a
  formal game-theoretic model showing that handicap equilibria could be
  evolutionarily stable.  See Grafen, A., ``Biological signals as
  handicaps,'' \textit{Journal of Theoretical Biology}, 144(4), 1990,
  pp.~517--546.}

Sosis and Alcorta~\cite{sosis2003} extended this logic from biology to
religion, arguing that costly rituals---painful initiations, dietary
restrictions, time-consuming liturgical practices---function as
\emph{credibility-enhancing displays} (CREDs) that solve the free-rider
problem in cooperative groups.  Religious communities face a persistent
threat from individuals who would enjoy the benefits of group membership
(mutual aid, shared childcare, collective defence) without bearing the
costs of commitment.  Costly rituals screen out free-riders by imposing
costs that only genuinely committed individuals are willing to pay.

Henrich's~(2009) cultural evolution framework~\cite{henrich2009}
situates costly signalling within a broader theory of cultural group
selection: groups that evolve effective costly-signalling mechanisms
outcompete groups that do not, because they achieve higher levels of
internal cooperation.\footnote{See also Boyd and Richerson (1985) on the
  coevolution of genes and culture, and Norenzayan (2013) on ``Big
  Gods'' as a cultural evolutionary response to the scaling problem of
  cooperation.}  Irons~(2001) independently developed a costly
signalling account of religious commitment, emphasising that religious
costs are ``hard to fake'' precisely because they involve genuine
sacrifice---time, resources, pain, or opportunity cost.  Bulbulia~(2004)
further argued that religious costs function as \emph{credibility
enhancers}: the greater the cost borne by a signaller, the more credible
the signal of underlying commitment.\footnote{Bulbulia, J., ``The
  cognitive and evolutionary psychology of religion,''
  \textit{Biology and Philosophy}, 19(5), 2004, pp.~655--686.}

The central contribution of Ideatic Depth Theory to this literature is
the identification of \textbf{depth with cost}.  Where Zahavi spoke of
metabolic cost, Spence of educational cost, and Sosis of ritual cost,
IDT provides a unified formal measure: the \emph{depth} of a signal in
an ideatic space.  This identification is not merely notational
convenience.  It allows us to derive, as formal theorems, the key
properties that costly signalling theorists have posited as assumptions:
subadditivity of combined signals, linear scaling of iterated rituals,
the single-crossing property that separates types, and the bounded
impact of signals on audiences of finite sophistication.

Zahavi's (1975) handicap principle~\cite{zahavi1975} and Sosis's (2003)
costly signalling theory of ritual~\cite{sosis2003} propose that honest
communication requires \emph{costly} signals.  The depth of an idea---its
structural complexity---is a natural cost metric: deeper signals require
more cognitive resources to produce and maintain.

In classical signalling theory (Spence~\cite{spence1973}), the cost of
education serves as a separating device between high-ability and
low-ability workers.  In our framework, the cost of a signal is its
\emph{depth}: the structural complexity that must be generated and
maintained.  A simple slogan has low cost; a multi-layered ritual
performance has high cost.

The key anthropological insight is that rituals involving genuine
sacrifice---scarification, potlatch, fasting---are \emph{honest depth
signals}: they cannot be cheaply faked because the depth (complexity) is
intrinsic to the act itself.

\medskip
\noindent\textbf{Chapter overview.}  We identify depth with signal cost
(\S\ref{sec:ch3:cost}), study iterated signals and ritual repetition
(\S\ref{sec:ch3:iteration}), formalise the single-crossing property
(\S\ref{sec:ch3:crossing}), bound interpretation cost
(\S\ref{sec:ch3:interp-cost}), and analyse the interaction between
resonance and cost (\S\ref{sec:ch3:resonance-cost}).

\subsection*{Historical Context: From Veblen to Zahavi}

The intuition that costliness signals quality has deep roots in the
social sciences, well predating Zahavi's biological formalisation.
Thorstein Veblen's \textit{Theory of the Leisure Class}~(1899)
introduced the concept of ``conspicuous consumption''---the display of
wealth through deliberately wasteful expenditure---as a mechanism for
signalling social status.\footnote{Veblen, T., \textit{The Theory of
  the Leisure Class: An Economic Study of Institutions}, Macmillan,
  1899.  Veblen's analysis of ``conspicuous waste'' anticipates the
  costly signalling framework by nearly a century.}  The industrial
magnate who commissions a mansion far beyond his residential needs, the
aristocrat who maintains a retinue of idle servants---these are,
in IDT terms, high-depth signals whose cost is the point.

Marcel Mauss's \textit{The Gift}~(1925) extended this logic to
non-market societies.  The potlatch ceremonies of the Pacific Northwest
involved competitive gift-giving and, in extreme cases, competitive
destruction of wealth: canoes broken, blankets burned, coppers thrown
into the sea.  Mauss recognised that these acts of apparent waste were
in fact \emph{investments in social obligation}: the depth of the
sacrifice created debts that bound communities
together.\footnote{Mauss, M., \textit{Essai sur le don}, L'Ann\'{e}e
  Sociologique, 1925.  English translation: \textit{The Gift}, trans.\
  W.~D.~Halls, Routledge, 1990.}

What Veblen and Mauss lacked was a formal framework capable of deriving
the \emph{consequences} of costly signalling from its
\emph{structure}.  Zahavi~(1975) provided the biological version of
such a framework; Spence~(1973) provided the economic version.  IDT
now provides a framework general enough to encompass both, grounding
cost in the combinatorial structure of ideatic depth.  The theorems that
follow make precise the properties that Veblen, Mauss, Zahavi, and
Sosis described informally: why cost is subadditive, why empty signals
carry no cost, and why high-quality types bear cost more
efficiently than low-quality types.

% ── §1. Depth as Signal Cost ─────────────────────────────────────────
\section{Depth as Signal Cost}
\label{sec:ch3:cost}

\begin{definition}[Signal Cost]
\label{def:signal-cost}
The \emph{cost} of signal~$s$ is its depth:
$\operatorname{signalCost}(s) = \depth(s)$.
In Zahavi's handicap principle, this is the metabolic, social, or
cognitive cost of producing the signal.
\end{definition}

\begin{theorem}[Void is Costless]
\label{thm:void-costless}
$\operatorname{signalCost}(\void) = 0$.
\end{theorem}

\begin{proof}[Proof sketch]
By $\depth(\void) = 0$.  See \texttt{void\_costless}.
\end{proof}

\begin{remark}
Silence is always free.  Every signalling system must contend with the
zero-cost outside option of saying nothing.
\end{remark}

\begin{theorem}[Cost Subadditivity]
\label{thm:cost-subadditive}
\[
  \operatorname{signalCost}(s_1 \compose s_2)
  \leq \operatorname{signalCost}(s_1) + \operatorname{signalCost}(s_2).
\]
\end{theorem}

\begin{proof}[Proof sketch]
By the depth-compose bound.  See \texttt{cost\_subadditive}.
\end{proof}

\begin{remark}
A two-part ritual (fasting + scarification) costs at most as much as the
parts separately.  Compound signals can be \emph{cheaper} than their
parts (subadditivity), but never more expensive.  This formalises the
``economies of ritual combination'' observed in initiation ceremonies
that bundle multiple ordeals.
\end{remark}

\begin{theorem}[Void Right-Compose Cost]
\label{thm:void-right-compose-cost}
$\operatorname{signalCost}(s \compose \void) = \operatorname{signalCost}(s)$.
\end{theorem}

\begin{proof}[Proof sketch]
By $\depth(s \compose \void) = \depth(s)$.
See \texttt{void\_compose\_cost}.
\end{proof}

\begin{theorem}[Void Left-Compose Cost]
\label{thm:void-left-compose-cost}
$\operatorname{signalCost}(\void \compose s) = \operatorname{signalCost}(s)$.
\end{theorem}

\begin{proof}[Proof sketch]
By $\depth(\void \compose s) = \depth(s)$.
See \texttt{void\_left\_compose\_cost}.
\end{proof}

\begin{remark}
You cannot fake costly signals by ``adding silence.''  Prefixing or
appending a dramatic pause to a sacrifice does not change its cost.
\end{remark}

\begin{example}[The Potlatch of the Pacific Northwest]
Among the Kwakwaka'wakw (formerly called Kwakiutl) peoples of the
Pacific Northwest coast, the potlatch ceremony served as a paradigmatic
costly signal.  Chiefs competed for social rank by giving away---and in
extreme cases, deliberately destroying---enormous quantities of wealth:
blankets, canoes, eulachon oil, and prized copper shields known as
\textit{t'lakwa}.  The depth of the signal was measured not in symbolic
gestures but in real material sacrifice.  A chief who destroyed a copper
worth thousands of blankets was producing a signal of such depth that no
rival of lesser means could replicate it.  In IDT terms,
$\operatorname{signalCost}(\text{potlatch}) = \depth(\text{potlatch})$,
and this depth was precisely what made the signal unfakeable: one cannot
pretend to destroy wealth one does not possess.\footnote{Boas, F.,
  \textit{Kwakiutl Ethnography}, ed.\ H.~Codere, University of Chicago
  Press, 1966.  See also Codere, H., ``Fighting with property: a study
  of Kwakiutl potlatching and warfare,'' American Ethnological Society
  Monograph~18, 1950.}
\end{example}

\begin{example}[Scarification Among the Nuer]
Evans-Pritchard's ethnography of the Nuer of southern Sudan documented a
striking initiation rite: adolescent boys received six deep horizontal
cuts across the forehead, from ear to ear, using a small knife.  The
cuts were made without anaesthetic, and the resulting scars---called
\textit{gaar}---were permanent, visible markers of adult male
status.\footnote{Evans-Pritchard, E.~E., \textit{The Nuer: A
  Description of the Modes of Livelihood and Political Institutions of
  a Nilotic People}, Clarendon Press, 1940.}  The signal satisfies the
formal properties of costly depth: the pain and risk of infection
constitute genuine cost that cannot be faked (no one can acquire
\textit{gaar} scars without enduring the cutting), and the permanence of
the signal means that its depth is \emph{preserved over time}---it
cannot be retracted or diminished.  By Theorem~\ref{thm:void-costless},
void signals (mere words or gestures) have zero cost; the Nuer
initiation achieves honesty precisely by requiring non-zero depth.
\end{example}

\subsection*{Discussion: Why Depth Cannot Be Faked}

The examples above illustrate a crucial structural property of depth as
a cost metric: depth is \emph{intrinsic} to the signal, not extrinsic.
A potlatch chief cannot fake the destruction of a copper because the
copper must actually be broken; a Nuer initiate cannot fake the
scarification because the scars must actually be cut.  This
intrinsicality is what distinguishes depth-as-cost from purely
conventional costs (such as monetary fees, which can be borrowed,
embezzled, or counterfeited).  In the formal framework, this property
is captured by the fact that $\depth(s)$ is determined by the
compositional structure of~$s$ itself, not by any external accounting
mechanism.  A signal's cost is its complexity, and complexity cannot be
counterfeited.

\begin{theorem}[Trivial Composition is Costless]
\label{thm:trivial-compose-costless}
If $\operatorname{signalCost}(s_1) = 0$ and
$\operatorname{signalCost}(s_2) = 0$, then
$\operatorname{signalCost}(s_1 \compose s_2) = 0$.
\end{theorem}

\begin{proof}[Proof sketch]
By the depth-zero-closed property of ideatic spaces.
See \texttt{trivial\_compose\_costless}.
\end{proof}

\begin{remark}
Banality composed with banality remains banal.  You cannot bootstrap
costly signalling from costless components---depth must come from
somewhere.
\end{remark}

% ── §2. Iterated Signals ─────────────────────────────────────────────
\section{Iterated Signals and Ritual Repetition}
\label{sec:ch3:iteration}

\begin{theorem}[Iteration Cost Bound]
\label{thm:iteration-cost}
For $n$-fold repetition of signal~$s$:
\[
  \operatorname{signalCost}(\operatorname{composeN}(s, n))
  \leq n \cdot \operatorname{signalCost}(s).
\]
\end{theorem}

\begin{proof}[Proof sketch]
By induction on~$n$, using subadditivity at each step.
See \texttt{iteration\_cost\_bound}.
\end{proof}

\begin{remark}
Ritual repetition (daily prayer, weekly sabbath, annual pilgrimage) has
linearly bounded cumulative cost.  A devotee who prays 5 times daily
pays at most $5 \times$ the cost of a single prayer.  This linear
scaling is why religious systems can demand iterative commitment without
unbounded cost.
\end{remark}

\begin{theorem}[Zero-Cost Iteration]
\label{thm:zero-cost-iteration}
$\operatorname{signalCost}(\operatorname{composeN}(\void, n)) = 0$
for all~$n$.
\end{theorem}

\begin{proof}[Proof sketch]
$\operatorname{composeN}(\void, n) = \void$ by the void-composeN
identity, so $\depth = 0$.  See \texttt{zero\_cost\_iteration}.
\end{proof}

\begin{remark}
Empty ritual (void composed with void, repeated endlessly) never
generates genuine cost.  ``Going through the motions'' with no internal
depth is formally costless.
\end{remark}

\begin{remark}[Connection to Religious Studies: The Five Pillars of Islam]
The five daily prayers (\textit{salat}) of Islam exemplify iterated
costly signalling in a religious context.  Each individual prayer has
bounded cost: the time and cognitive effort required for the ablution
(\textit{wudu}), the physical prostrations (\textit{ruku} and
\textit{sujud}), and the recitation of Qur'anic passages.  By
Theorem~\ref{thm:iteration-cost}, the total annual cost of
\textit{salat} is linearly bounded by the number of iterations times
the cost of a single prayer.  What makes this system particularly
elegant from a signalling perspective is that the per-iteration cost is
calibrated to be bearable---Islam explicitly forbids making worship
unduly burdensome (\textit{la yukallifu Allahu nafsan illa wus'aha},
Qur'an 2:286)---while the cumulative commitment over a lifetime is
substantial enough to separate genuine believers from
opportunists.\footnote{Cf.\ Sosis, R., ``Why aren't we all
  Hutterites? Costly signaling theory and religious behavior,''
  \textit{Human Nature}, 14(2), 2003, pp.~91--127.}
\end{remark}

\begin{example}[Daily Prayer as Iterated Signal]
A Muslim who performs \textit{salat} five times daily for fifty years
executes approximately $5 \times 365 \times 50 = 91{,}250$ iterations
of the prayer signal.  By Theorem~\ref{thm:iteration-cost}, the total
signal cost over a lifetime is bounded above by
$91{,}250 \times \operatorname{signalCost}(\text{single prayer})$.
This linear bound has an important implication: it means that the
cumulative cost of religious observance grows \emph{predictably}, not
explosively.  A community can calibrate its ritual demands to be
sustainable across generations, which is exactly what we observe in
the major world religions.  By contrast, if signal cost were
superadditive (if each repetition made the next one more costly), no
religious system could sustain iterative practice over decades.
\end{example}

\begin{theorem}[Single Repetition Cost]
\label{thm:single-rep-cost}
$\operatorname{signalCost}(\operatorname{composeN}(s, 1))
 \leq \operatorname{signalCost}(s)$.
\end{theorem}

\begin{proof}[Proof sketch]
Specialise Theorem~\ref{thm:iteration-cost} with $n = 1$.
See \texttt{single\_rep\_cost}.
\end{proof}

% ── §3. The Single-Crossing Property ─────────────────────────────────
\section{The Single-Crossing Property}
\label{sec:ch3:crossing}

\begin{definition}[Costly Signal Model]
\label{def:costly-signal-model}
A \emph{costly signal model} consists of two cost functions:
\begin{itemize}
  \item $c_H : \IS \to \mathbb{N}$ (high-quality type), with
    $c_H(s) \leq \depth(s)$;
  \item $c_L : \IS \to \mathbb{N}$ (low-quality type), with
    $\depth(s) \leq c_L(s)$.
\end{itemize}
The high type can produce depth efficiently; the low type cannot.
\end{definition}

\begin{lstlisting}[caption={CostlySignalModel in Lean~4},label=lst:costly-signal]
structure CostlySignalModel (I : Type*) [IdeaticSpace I] where
  highTypeCost : I -> Nat
  lowTypeCost : I -> Nat
  high_efficient : forall s, highTypeCost s <= depth s
  low_costly : forall s, depth s <= lowTypeCost s
\end{lstlisting}

\begin{theorem}[High Type Cost Advantage]
\label{thm:high-type-advantage}
For any costly signal model and any signal~$s$:
$c_H(s) \leq c_L(s)$.
\end{theorem}

\begin{proof}[Proof sketch]
Chain the two constraints: $c_H(s) \leq \depth(s) \leq c_L(s)$.
See \texttt{high\_type\_advantage}.
\end{proof}

\begin{remark}
In Sosis's theory, genuine believers find costly rituals less onerous
than hypocrites.  The devout faster suffers less than the faker because
their internal ``type'' (depth of belief) makes the signal cheaper to
produce.  This is the formal foundation of ritual honesty.
\end{remark}

\begin{theorem}[Void is Free for All Types]
\label{thm:void-free-all}
$c_H(\void) \leq 0$ (and hence $c_H(\void) = 0$).
\end{theorem}

\begin{proof}[Proof sketch]
$c_H(\void) \leq \depth(\void) = 0$.
See \texttt{void\_free\_all\_types}.
\end{proof}

\begin{remark}
The babbling option is costless for all types.  Separation in costly
signalling must come from \emph{non-void} signals---only signals with
genuine depth can separate types.
\end{remark}

\subsection*{Discussion: The Ecology of Costly Signalling}

Why do different societies calibrate signal costs so differently?  The
Yanomam\"{o} demand chest-pounding duels; the Amish demand plain dress
and technological renunciation; the Benedictines demand a lifetime of
prayer, labour, and obedience.  From the perspective of IDT, these
are all costly signals, but their \emph{depth profiles} differ
dramatically.  The key insight from Sosis's comparative work is that
the optimal signal cost depends on the cooperative ecology of the
group: groups facing greater free-rider threats should evolve costlier
signals, and groups whose cooperation yields greater per-capita
benefits can afford to impose higher costs without losing members.

This prediction has been empirically tested.  Sosis and
Bressler~(2003) analysed the longevity of 277 communes founded in the
nineteenth-century United States, comparing religious communes (which
imposed costly ritual requirements) with secular communes (which
typically did not).\footnote{Sosis, R.\ and Bressler, E.~R.,
  ``Cooperation and commune longevity: a test of the costly signaling
  theory of religion,'' \textit{Cross-Cultural Research}, 37(2), 2003,
  pp.~211--239.}  They found that religious communes survived
significantly longer than secular ones, and that within the religious
category, communes imposing more costly requirements (dietary
restrictions, distinctive dress, behavioural taboos) survived longer
still.  Each additional costly requirement increased the predicted
lifespan of the commune.

\begin{example}[Kibbutzim and Religious Communes]
The contrast between religious and secular kibbutzim in
twentieth-century Israel provides a natural experiment in costly
signalling.  Religious kibbutzim required members to observe Shabbat,
keep kosher, and participate in daily prayer---all costly signals in
the IDT sense, with non-trivial depth.  Secular kibbutzim, by
contrast, relied primarily on ideological commitment, a signal whose
depth is harder to verify externally.  Consistent with costly
signalling theory, religious kibbutzim exhibited lower defection rates
and greater economic cooperation.\footnote{Sosis, R.\ and Ruffle,
  B.~J., ``Religious ritual and cooperation: testing for a
  relationship on Israeli religious and secular kibbutzim,''
  \textit{Current Anthropology}, 44(5), 2003, pp.~713--722.}
By Theorem~\ref{thm:high-type-advantage}, genuinely committed members
of religious kibbutzim bore these costs more efficiently than
uncommitted members, producing exactly the type-separation that
sustains long-term cooperation.
\end{example}

% ── §4. Interpretation Cost ──────────────────────────────────────────
\section{Interpretation Cost and Audience Effects}
\label{sec:ch3:interp-cost}

\begin{theorem}[Total Interpretation Cost Bound]
\label{thm:total-interp-cost}
\[
  \operatorname{depthSum}(\operatorname{interpret}(\Rep, s))
  \leq \operatorname{depthSum}(\Rep) + |\Rep| \cdot \operatorname{signalCost}(s).
\]
\end{theorem}

\begin{proof}[Proof sketch]
By the Total Hermeneutic Bound (Theorem~\ref{thm:total-hermeneutic-bound}),
noting that $\operatorname{signalCost}(s) = \depth(s)$.
See \texttt{total\_interpretation\_cost}.
\end{proof}

\begin{remark}
The ``impact'' of a costly signal on an audience is bounded by the
signal's cost times the audience size plus the audience's existing
complexity.  A potlatch (costly gift-giving) impresses proportionally to
its cost $\times$ audience.
\end{remark}

\begin{theorem}[Costless Signal Preserves Total Depth]
\label{thm:costless-preserves}
If $\operatorname{signalCost}(s) = 0$, then
$\operatorname{depthSum}(\operatorname{interpret}(\Rep, s))
 \leq \operatorname{depthSum}(\Rep)$.
\end{theorem}

\begin{proof}[Proof sketch]
Set $\depth(s) = 0$ in the hermeneutic bound to eliminate the linear
term.  See \texttt{costless\_preserves\_depth}.
\end{proof}

\begin{theorem}[Atomic Signal Bounded Cost]
\label{thm:atomic-cost}
If $s$ is atomic, then $\operatorname{signalCost}(s) \leq 1$.
\end{theorem}

\begin{proof}[Proof sketch]
By $\depth(s) \leq 1$ for atomic ideas.
See \texttt{atomic\_signal\_cost}.
\end{proof}

\begin{theorem}[Compound Atomic Cost]
\label{thm:compound-atomic-cost}
If $s_1, s_2$ are both atomic, then
$\operatorname{signalCost}(s_1 \compose s_2) \leq 2$.
\end{theorem}

\begin{proof}[Proof sketch]
By $\depth(s_1 \compose s_2) \leq \depth(s_1) + \depth(s_2) \leq 2$.
See \texttt{compound\_atomic\_cost}.
\end{proof}

\begin{remark}
The simplest possible compound ritual (two atomic acts) has strictly
bounded cost---complex enough to potentially separate types but simple
enough to be tractable.
\end{remark}

\begin{remark}[Connection to Evolutionary Biology]
The compound atomic cost theorem
(Theorem~\ref{thm:compound-atomic-cost}) has a direct biological
analogue in multi-component mating displays.  The peacock's courtship
involves both a visual component (the iridescent tail fan, with its
elaborate eyespot pattern) and an auditory component (a distinctive
rattling call produced by vibrating the tail feathers).  Each component
is approximately atomic in the IDT sense: a single, irreducible costly
act.  The total display cost is bounded by the sum of the component
costs, exactly as the theorem predicts.  Crucially, the two components
are not redundant---they are received through different sensory channels
and serve complementary verification functions.  The visual component
signals genetic quality (symmetry, pigmentation); the auditory component
signals current condition (vigour, muscular control).  The compound
signal is more informative than either component alone, while remaining
within the tractable cost bound of~$2$.\footnote{Loyau, A.\ et al.,
  ``Iridescent structurally based coloration of eyespots correlates
  with mating success in the peacock,'' \textit{Behavioral Ecology},
  18(6), 2007, pp.~1123--1131.}
\end{remark}

\begin{example}[Marriage Ceremonies as Compound Signals]
A traditional Hindu wedding (\textit{vivaha}) provides an instructive
example of compound costly signalling.  The ceremony involves multiple
individually costly acts: the exchange of garlands (\textit{jaimala}),
signalling mutual acceptance; the sacred fire ritual (\textit{agni
parikrama}), in which the couple circles the fire seven times while
reciting vows; the application of vermilion (\textit{sindoor}) to the
bride's hair parting; and the tying of the sacred thread
(\textit{mangalsutra}).  Each act is approximately atomic---a single
ritual gesture with bounded depth---and the total ceremony has bounded
compound cost by Theorem~\ref{thm:compound-atomic-cost} applied
iteratively.  The multi-component structure serves a dual purpose:
it increases the total signal cost (making the commitment harder to
fake) while distributing the cost across distinct symbolic registers
(visual, verbal, kinesthetic, material), thereby making the signal
legible to a diverse audience of witnesses with different interpretive
repertoires.\footnote{Pandey, R., \textit{Hindu Sa\d{m}sk\={a}ras:
  Socio-Religious Study of the Hindu Sacraments}, Motilal Banarsidass,
  1969.}
\end{example}

% ── §5. Resonance and Cost ───────────────────────────────────────────
\section{Resonance and Cost Interaction}
\label{sec:ch3:resonance-cost}

\begin{theorem}[Resonant Signals Have Comparable Outcomes]
\label{thm:resonant-comparable}
If $s_1 \res s_2$, then for each $r \in \Rep$,
$r \compose s_1 \res r \compose s_2$.
\end{theorem}

\begin{proof}[Proof sketch]
By resonance-compose-left applied to each repertoire element.
See \texttt{resonant\_signals\_comparable}.
\end{proof}

\begin{remark}
Two rituals that ``resonate'' produce similar effects on any audience,
even if one costs more than the other.  This is why cheap imitations of
costly rituals can sometimes ``work''---if they resonate with the
original, the audience's interpretations resonate.
\end{remark}

\begin{theorem}[Self-Resonant Interpretation]
\label{thm:self-resonant-interp}
For any repertoire~$\Rep$ and signal~$s$,
$\operatorname{interpret}(\Rep, s)$ is pairwise self-resonant.
\end{theorem}

\begin{proof}[Proof sketch]
Apply Theorem~\ref{thm:resonant-comparable} with $s_1 = s_2 = s$,
using reflexivity of resonance.
See \texttt{self\_resonant\_interpretation}.
\end{proof}

\begin{theorem}[Bounded Repertoire, Bounded Outcome]
\label{thm:bounded-repertoire}
If $\depth(r) \leq d$ for all $r \in \Rep$, then each interpretation
$x \in \operatorname{interpret}(\Rep, s)$ satisfies
$\depth(x) \leq d + \operatorname{signalCost}(s)$.
\end{theorem}

\begin{proof}[Proof sketch]
For each $r \in \Rep$,
$\depth(r \compose s) \leq \depth(r) + \depth(s) \leq d + \depth(s)$.
See \texttt{bounded\_repertoire\_bounded\_outcome}.
\end{proof}

\begin{remark}
Receivers with bounded sophistication cannot produce arbitrarily deep
interpretations.  Even the most costly signal only adds its own cost to
each receiver's maximum capacity.
\end{remark}

\begin{theorem}[Interpretation Filtration]
\label{thm:interp-filtration}
If $\Rep \subseteq \mathcal{F}_n$ (the depth-$n$ filtration), then
$\operatorname{interpret}(\Rep, s) \subseteq
 \mathcal{F}_{n + \operatorname{signalCost}(s)}$.
\end{theorem}

\begin{proof}[Proof sketch]
Each $r \in \Rep$ has $\depth(r) \leq n$, so
$\depth(r \compose s) \leq n + \depth(s)$.
See \texttt{interpretation\_filtration}.
\end{proof}

\begin{remark}
Costly signals push interpretations up the depth hierarchy by at most
their cost.  A depth-3 ritual received by depth-5 minds produces at most
depth-8 interpretations.  This is the formal ``budget'' of cultural
transformation through signalling.
\end{remark}

\begin{remark}[Open Questions in Costly Signalling]
Several questions remain open in the formal theory of costly depth
signals, and we flag them here as directions for future work.
First, can the \emph{single-crossing property}
(Theorem~\ref{thm:high-type-advantage}) be weakened?  Our current
formulation assumes a binary type space (high vs.\ low), but real
signalling environments feature a continuum of types.  Extending the
single-crossing result to a richer type space would require a graded
notion of cost efficiency, perhaps indexed by a continuous depth
parameter.  Second, what happens when signal cost is
\emph{stochastic}?  In practice, the cost of a ritual performance
varies across occasions: a fast during summer is more costly than a
fast during winter; a pilgrimage in wartime is more costly than in
peacetime.  A stochastic extension of IDT would model
$\operatorname{signalCost}(s)$ as a random variable and derive
probabilistic analogues of the theorems in this chapter.  Third, the
interaction between \emph{signal cost and audience size} deserves
further exploration: Theorem~\ref{thm:total-interp-cost} gives a
linear bound, but empirical evidence suggests that audience effects may
exhibit diminishing returns (a potlatch before 1000 witnesses is not
ten times more effective than one before 100).
\end{remark}

\begin{example}[The Olympic Games as Multi-Level Costly Signal]
The modern Olympic Games function as a costly signal at multiple
levels simultaneously, illustrating the compositional structure that
IDT formalises.  At the \emph{national} level, the host country bears
enormous costs: infrastructure construction, security, and
opportunity cost of diverted public funds.  These costs signal
national prosperity and organisational capacity.  At the
\emph{individual} level, each athlete bears years of training cost:
the depth of their athletic signal is measured in thousands of hours
of deliberate practice, dietary discipline, and foregone alternative
careers.  At the \emph{spectator} level, attendees bear costs of
travel, tickets, and time---signals of allegiance to their national
team.  By the subadditivity theorem
(Theorem~\ref{thm:cost-subadditive}), the total cost of the Olympic
signal is bounded by the sum of these component costs, and by the
interpretation filtration theorem
(Theorem~\ref{thm:interp-filtration}), the cultural impact of the
Games on any audience is bounded by the audience's existing
sophistication plus the signal's depth.\footnote{For a sociological
  analysis of the Olympics as costly signalling, see Roche, M.,
  \textit{Mega-Events and Modernity: Olympics and Expos in the Growth
  of Global Culture}, Routledge, 2000.}
\end{example}

% ── Summary ──────────────────────────────────────────────────────────
\section*{Chapter Summary}

This chapter connected IDT depth to signal cost theory.  The
identification $\operatorname{signalCost}(s) = \depth(s)$ allowed us to
prove that void is costless (Theorem~\ref{thm:void-costless}), cost is
subadditive (Theorem~\ref{thm:cost-subadditive}), and ritual repetition
has linearly bounded cumulative cost
(Theorem~\ref{thm:iteration-cost}).  The \emph{single-crossing property}
(Theorem~\ref{thm:high-type-advantage}) establishes that high-quality
types pay less for depth, providing the formal foundation of Zahavi's
handicap principle and Sosis's costly signalling theory of ritual.
Resonant signals produce comparable outcomes regardless of cost
difference (Theorem~\ref{thm:resonant-comparable}), and the
\emph{interpretation filtration} theorem
(Theorem~\ref{thm:interp-filtration}) shows that costly signals push
interpretations up the depth hierarchy by exactly their cost.

All eighteen theorems are verified in
\texttt{FormalAnthropology/IDT\_Signal\_Ch03.lean}.
